% Context from chapters-tex/approximation.tex; for mathematical reference, not compilation.
\section{Approximation of Cartoon Images}

The isotropic, compactly supported wavelets considered here occupy squares that do not follow the smooth contours of geometric images~\eqref{eq-cartoon-model}. These contours have more structure than the finite-length level sets allowed by the bounded-variation model~\eqref{eq-model-tv}.

                                        
\subsection{Wavelet Approximation of Cartoon Images}

The bound~\eqref{eq-approximation-error-wav-2d} gives a wavelet approximation rate of $O(M^{-1})$ for cartoon images~\eqref{eq-cartoon-model}. Even the simple indicator $f=1_\Om$, where $\Om$ is a disk, can attain this slow rate: isotropic wavelets do not efficiently follow its curved boundary. The smoothed cartoon model~\eqref{eq-cartoon-model-smooth} can behave similarly when the coefficient budget $M$ is too small to resolve the blur.

Figure \ref{fig-cartoon-wav} shows that many large coefficients are located near edge curves, and retaining only a small number leads to a poor approximation with visually unpleasant artifacts.


\myfigure{
\tabtrois{
\image{approximation}{.32}{cartoon-wav-image}&
\image{approximation}{.32}{cartoon-wav-approx}&
\image{approximation}{.32}{cartoon-wav-coefs}\\
$f$ & $f_M$ & $\{\dotp{f}{\psi_{j,n}^\om}\}$
}
}{ 
	Wavelet approximation of a cartoon image.   
}{fig-cartoon-wav}

                                        
\subsection{Finite Element Approximation}

An adaptive triangulation can improve the approximation by using elongated triangles aligned with the edges.
Figure \ref{fig-triangulation-image} shows an example of such a triangulation.

\myfigure{
\image{approximation}{.3}{triangulation-image}
\image{approximation}{.3}{triangulation-triangles}
}{ 
	Left: cartoon image, right: adaptive triangulation.   
}{fig-triangulation-image}

Select $M$ points in the image domain $[0,1]^2$ and connect them into a triangulation. Define $\tilde f_M$ by piecewise affine interpolation on the resulting triangles.

\myfigure{
\image{approximation}{.3}{triangulation-aniso}\quad
\image{approximation}{.3}{triangulation-iso}
}{ 
	Anisotropic triangles near an edge (left) and approximately isotropic triangles away from edges (right).  
}{fig-triangulation-anisoiso}

Figure \ref{fig-triangulation-anisoiso} illustrates a construction for $\Cdeux$ cartoon images~\eqref{eq-cartoon-model} with $\al=2$. In smooth regions, use $\approx M/2$ nearly equilateral triangles of width $\approx M^{-1/2}$. Near the edges, exploit the $\Cdeux$ contour regularity by placing $\approx M/2$ elongated triangles of length $\approx M^{-1}$ along the contours and width $\approx M^{-2}$ across them. Such a triangulation yields the error bound
\eql{\label{eq-error-triangulation}
	\norm{f-f_M}^2 = O(M^{-2}),
}
which improves over the wavelet approximation error decay \eqref{eq-approximation-error-wav-2d}. 

Implementing this construction requires estimating the edges. This is particularly difficult for blurred cartoon images with an unknown smoothing kernel $h$.

Practical methods use heuristics or greedy rules to select sample points and triangulate them from discrete or noisy data.
Figure~\ref{fig-triangulation-peppers} illustrates compression with one such greedy construction.

\myfigure{
\tabtrois{
\image{approximation}{.32}{triangulation-peppers}&
\image{approximation}{.32}{triangulation-peppers-triangles}&
\image{approximation}{.32}{triangulation-peppers-jpeg2k}\\
Adapted triangulation & $\tilde f_M$ & JPEG-2000
}
}{ 
	Comparison of adaptive triangulation and JPEG-2000, with the same number of bits.   
}{fig-triangulation-peppers}

                                        
\subsection{Curvelet Approximation}

A fixed family of oriented, anisotropic atoms offers an alternative to adaptive triangulations. Cand\`es and Donoho introduced the curvelet frame for this purpose~\cite{candes-tight-frame}.

  
\paragraph{Curvelets.}

At fine scales $j\leq0$, the geometry can be illustrated by applying parabolic scaling to a horizontally oriented mother function $c$,
\eq{
	c_{2^j}(x_1,x_2) \approx 2^{-3j/4} c(2^{-j/2} x_1, 2^{-j} x_2),
} 
followed by translation and rotation,
\eq{
	c_{2^j,u}^\th(x_1,x_2) = c_{2^j}(R_{-\th}(x-u))
}
where $R_\th$ is the rotation of angle $\th$.

The atom $c_{2^j,u}^\th$ is localized near $u$ and oriented at angle $\th$. Its dimensions obey ``width $\approx$ length$^2$'', the same relation used for adaptive triangles near edges. This scaling matches the contour geometry of the cartoon model~\eqref{eq-cartoon-model} with $\al=2$.

\myfigure{
\image{approximation}{.36}{curvelet-spacial}
\image{approximation}{.4}{curvelet-fourier}
}{ 
	A curvelet $c_m$ in space (left) and its localization in the Fourier domain (right). 
}{fig-curvelet}

Figure \ref{fig-curvelet} illustrates this localization in space and frequency.

  
\paragraph{Parameter discretization.}

To build an image representation, one needs to sample the position $u$ and orientation $\theta$.
The angular sampling interval depends on the scale:
\eq{
	\foralls 0 \leq k < 2^{-\lceil j/2 \rceil+2}, \quad
	\th_{k}^{(j)} = k \pi 2^{ \lceil j/2 \rceil -1 }  
}
and the spatial grid depends on both scale and orientation:
\eq{
	\foralls m=(m_1,m_2) \in \ZZ^2, \quad
	u_m^{(j,\th)} = R_\th( 2^{j/2} m_1, 2^j m_2 ).
}
Figure \ref{fig-curvelet-discretization} shows this sampling pattern.

\myfigure{
\image{approximation}{.8}{curvelet-discretization}
}{ 
	Sampling pattern for the curvelet positions.  
}{fig-curvelet-discretization}

  
\paragraph{Curvelet tight frame.}

The scaling and sampling formulas above describe the geometry o